\documentclass[11pt,a4paper]{article}
\usepackage[utf8]{inputenc}
\usepackage[T1]{fontenc}
\usepackage[magyar]{babel}
\usepackage{amsmath,amssymb}
\usepackage[margin=2.5cm]{geometry}
\usepackage{tcolorbox}
\usepackage{xcolor}
\usepackage{tikz}

% Színek definíciója a barátságos megjelenéshez
\definecolor{myprimary}{RGB}{41, 128, 185}
\definecolor{mysecondary}{RGB}{236, 240, 241}
\definecolor{myaccent}{RGB}{211, 84, 0}
\definecolor{myboxbg}{RGB}{249, 249, 249}

\title{Ismétlés mélküli variáció}
\author{A Te Kedves Magántanárod}
\date{}

%%% Megjelenítéshez automatikusan kiegészített csomagok %%%
\usepackage{amsfonts, amsthm}
\usepackage{pgfplots}
\pgfplotsset{compat=1.18}
\usepackage{graphicx}
\usepackage{enumitem}
\usepackage{booktabs}
\usepackage{float}
\usepackage{caption}
\usepackage{microtype}

\usepackage{hyperref}

\begin{document}

\maketitle

\begin{tcolorbox}[colback=mysecondary,colframe=myprimary,arc=4pt,title=\textbf{Üdvözöllek, Kedves Diákum!}]
Ne aggódj, ha a kombinatorika elsőre bonyolultnak tűnik! Sokan megijednek tőle, pedig ez lényegében nem más, mint logikus pakolgatás és számlálás. Vedd úgy, hogy most is csak egy kávé vagy tea mellett beszélgetünk a feladatokról. Vágjunk is bele!
\end{tcolorbox}

\section{Lélekmegnyugtató, lényegretörő nyitás}
Az \textbf{ismétlés nélküli variáció} nem más, mint amikor egy nagyobb készletből kiválasztunk néhány elemet, \textbf{számít a sorrend}, és \textbf{egy elemet csak egyszer} használhatunk fel. Képzeld el, hogy egy ötfős baráti társaságból kell kiválasztanod hármat, akik a dobogóra állhatnak (arany-, ezüst- és bronzéremért). Hányféleképpen alakulhat a dobogósok sorrendje? Na, ezt nevezzük ismétlés nélküli variációnak!

\section{Hogyan működik?}
Képzeljük el, hogy egy futóverseny döntőjében 8 versenyző indul ($n = 8$), és nekünk az első 3 helyezettet (arany-, ezüst-, bronzérem) kell kiosztanunk ($k = 3$). 

\begin{itemize}
    \item Hányféleképpen oszthatjuk ki az \textbf{aranyérmet}? Bármelyik 8 versenyző kaphatja. Ez $8$ lehetőség.
    \item Miután az aranyat kiosztottuk, hány ember maradt az \textbf{ezüstéremre}? Már csak $7$ versenyző. (Ezért \emph{ismétlés nélküli}, hiszen egy ember nem kaphat két érmet!)
    \item Végül, hányan maradtak a \textbf{bronzéremre}? Mivel kettőt már díjaztunk, maradt $6$ ember.
\end{itemize}

A kombinatorikai alapelv szerint minden egyes aranyérmes lehetőséghöz $7$ különböző ezüstérmes helyzet tartozik, és az így kapott párok mindegyikéhez még $6$ bronzos leosztás kapcsolódik. Tehát az összes lehetőség:
\[ 8 \cdot 7 \cdot 6 = 336 \]
Látod? Nem varázslat, csak tiszta logika!

\section{Általános képlet és speciális esetek}
Ha $n$ darab elemünk van, és ezek közül választunk ki sorrendben $k$ darabet (ahol $k \le n$), akkor a lehetőségek számát az alábbi \textbf{variációs képlet} adja meg:

\[ V_n^k = n \cdot (n-1) \cdot (n-2) \cdot \dots \cdot (n - k + 1) \]

Ezt a faktoriálisok segítségével sokkal elegánsabban is felírhatjuk:
\[ V_n^k = \frac{n!}{(n-k)!} \]

\begin{tcolorbox}[colback=myboxbg,colframe=myaccent,title=\textbf{Vigyázz a csapdákra!}]
\begin{itemize}
    \item \textbf{Sorrend:} Mindig kérdezd meg magadtól: számít-e, hogy milyen sorrendben húzzuk ki az elemeket? (Itt \textbf{igen}, a 123 nem ugyanaz, mint a 321).
    \item \textbf{Visszatevés:} Figyelj arra, hogy a szöveg említi-e, hogy visszatesszük az elemet. Ha visszatesszük, akkor az már \emph{ismétléses} variáció (ott $n^k$ a képlet), nem pedig ez!
\end{itemize}
\end{tcolorbox}

\section{Dolgozat- / Vizsga-gyorsteszt (Döntési fa)}
Amikor olvasod a feladatot a dolgozatban, tedd fel magadnak ezt a 3 kérdést sorban:
\begin{enumerate}
    \item \textbf{A teljes készletből csak egy részét választom ki?} (Ha mindet sorba rendezem, az csak egy sima $n!$ permutáció).
    \item \textbf{Számít a sorrend?} (Ha a sorrend nem számít, például egymaroknyi kártyánál vagy csapatenkénti összegyűjtésnél, akkor ez \emph{kombináció} lesz, nem variáció!).
    \item \textbf{Visszakerül az elem a kalapba húzás után?} (Ha igen $\rightarrow$ ismétléses variáció. Ha nem $\rightarrow$ \textbf{ez az ismétlés nélküli variáció}, amiről most tanulunk!).
\end{enumerate}

\section{3 tipikus és trükkös feladat levezetése}

\subsection*{1. feladat: A klasszikus jelszó-kód}
\textbf{Feladat:} Hány olyan 4-jegyű szám létezik, amelyben a számjegyek különbözőek, és a 0-val sem kezdődhetnek?

\noindent\textbf{Gondolatmenet:} 
Mit látunk? Van 10 számjegyünk ($0$-tól $9$-ig), és ki kell választanunk belőle 4-et úgy, hogy számít a sorrend (mert a 1234 nem ugyanaz, mint a 4321) és nem ismétlődhetnek. 
Itt egy pici csapda van: a szám nem kezdődhet 0-val!

\noindent\textbf{Megoldás lépésről lépésre:}
\begin{enumerate}
    $1.$ Válasszuk ki először az \textbf{első jegyet}. Mivel a 0 nem lehet, csak $9$ féle választásunk maradt ($1$-től $9$-ig).
    $2.$ Most jön a \textbf{második jegy}. Itt már szerepelhet a 0 is, de az első helyre már elhasználtunk egy számjegyet, így a maradék $10 - 1 = 9$ számjegy közül választhatunk.
    $3.$ A \textbf{harmadik helyre} $8$ féle számjegy jöhet.
    $4.$ A \textbf{negyedik helyre} pedig $7$ féle számjegy maradt.
\end{enumerate}
Az összes lehetséges szám: 
\[ 9 \cdot 9 \cdot 8 \cdot 7 = 4536 \]
\emph{Hol rontják el gyakran?} Sok diák elfelejti a 0-ás megkötést, és az első helyre is 10-et ír, vagy elfelejti, hogy a 0 miatt a második helyre mégiscsak visszajön egy opció.

\subsection*{2. feladat: Az elnök és az alelnök esete}
\textbf{Feladat:} Egy 12 fős baráti társaságból megválasztunk egy elnököt, egy alelnököt és egy pénztárost. Hányféle kimenetele lehet ennek a választásnak?

\noindent\textbf{Gondolatmenet:} 
Van 12 emberünk ($n=12$), és kiválasztunk 3-at ($k=3$). Mivel a pozíciók különbözőek (az elnöki szék más, mint a pénztárosé), a \textbf{sorrend szigorúan számít}. Egy ember nem kaphat két különböző főtisztséget egyszerre (ismétlés nincs). Ez pont az ismétlés nélküli variáció tankönyvi esete!

\noindent\textbf{Megoldás lépésről lépésre:}
Használhatjuk a képletünket:
\[ V_{12}^3 = \frac{12!}{(12-3)!} = \frac{12!}{9!} \]
Írjuk ezt ki szorzat alakban, és lássuk, hogyan egyszerűsödik:
\[ \frac{12 \cdot 11 \cdot 10 \cdot 9!}{9!} = 12 \cdot 11 \cdot 10 = 1320 \]
\textbf{Eredmény:} 1320-féleképpen alakulhat a vezetőség.

\begin{tcolorbox}[colback=mysecondary,colframe=myprimary,arc=4pt,title=\textbf{Záró gondolat}]
Látod? Ha lépésről lépésre haladsz, és felteszed a helyes kérdéseket, a kombinatorika egy kifejezetten szórakoztató logikai játék. Gyakorolj bátran, és hamarosan rutinná válik!
\end{tcolorbox}

\end{document}